\chapter{The Unnameable}
% OUTLINE Ch4 "The Unnameable (Tarski 1933)". Laws: both inherited
% (labeled intuitions adds-no-claim); non-funge. CITATION: Tarski 1933
% (undefinability of truth; "Der Wahrheitsbegriff..."); the semantic
% room as the syntactic room's twin; the object/meta hierarchy as the
% honest exit (the tower the series' readers know the feel of). FENCE:
% undefinability to FORMAL languages; no leap to natural language / the
% world. Zero device appearances. Gate: Kodo.

\begin{intuition}
A dictionary defines its words using other words, all of them in the same
dictionary, and mostly this works. But ask the dictionary to define one
particular word --- \emph{true} --- as a predicate that decides, for every
sentence the dictionary can write, whether that sentence is true, and a
familiar visitor arrives. Consider the sentence that says of itself, using
the dictionary's own new word, \emph{this sentence is not true}. If it is
true, then what it says holds, so it is not true; if it is not true, then
what it says holds, so it is true. The dictionary cannot house the word
without housing the contradiction. The feeling to carry: a language rich
enough to talk about its own sentences cannot also contain its own
faithful word for their truth --- not because truth is mysterious, but
because the word, added, would let the diagonal speak.
\end{intuition}

This is the second great room, and it is the first's twin. Where G\"odel
let arithmetic speak of its own \emph{provability} and found a sentence
the house could not prove, Alfred Tarski, in his 1933 monograph on the
concept of truth in formalized languages, let a language reach for its own
\emph{truth} and found the word could not live in the house at all. The
two results are the same gesture aimed one notch over: G\"odel's target
was the proof relation, a matter of syntax --- which strings the rules can
derive; Tarski's target is the truth relation, a matter of semantics ---
which sentences actually hold. Provability is a room inside the house;
truth turns out to be a door the house cannot have.

The argument is the liar sentence made rigorous, and it runs on the same
machinery as Chapter 3. Suppose a formal language could define, in its own
terms, a predicate ``true-of'' that correctly applied to exactly the
numbers coding its true sentences. Arithmetization --- G\"odel's bridge,
reused --- lets the language build, by the diagonal lemma, a sentence that
asserts its own untruth: provably equivalent, inside the language, to
``I am not true.'' Now the two branches close on each other with no exit,
the liar with a certificate, and the only assumption that can be dropped
to escape the contradiction is the one that started it: no such predicate
is definable within the language. Truth, for a sufficiently rich formal
language, is not expressible in that language. That is Tarski's theorem of
the undefinability of truth, and it is why the diagonal, here, does not
build a room inside the house but proves a door cannot be cut in its wall.

The honest exit is a staircase, and the reader of this series has climbed
its like before. Truth for a language can be defined --- rigorously,
completely --- but only in a \emph{richer} language, one step up, that can
speak of the first language's sentences from outside. Call the first the
object language and the second the metalanguage; the metalanguage has its
own truth, undefinable in \emph{it}, definable only one step higher again.
Truth is not destroyed by Tarski's theorem; it is \emph{stratified} ---
each language's truth living one floor above it, in an ascending tower
with no top floor. The reader who has felt the earlier volumes' towers ---
the descriptions describing descriptions, the frames read from frames ---
will recognize the shape: not a wall against truth, but a refusal to let
any single level hold its own. The room the diagonal makes here is the
gap between a floor and the floor that can speak its truth, and the gap
never closes, only climbs.

\begin{intuition}
The picture: a language is a room with a window, and through the window it
can see and name everything outside --- except the room it is standing in.
To describe this room truthfully you must step into the hallway, which is
another room, with its own blind spot. There is no room with a window onto
itself. The tower of metalanguages is that hallway, going up forever, each
floor naming the truth of the one below and blind to its own.
\end{intuition}

Two honesty notes, and the first is the fence this chapter must hold.

Tarski's theorem is about \emph{formal languages} of a certain strength,
and its exit --- the object/meta stratification --- is a fact about such
languages, proved. It is tempting, and common, to read it as a verdict on
\emph{natural} language: that English, say, cannot coherently contain its
own word ``true,'' or that ordinary talk of truth is incoherent, or that
Tarski proved truth to be relative. The theorem supports none of this.
Natural languages plainly do contain the word ``true'' and use it about
their own sentences daily; whether they do so consistently, and how, is a
question in the philosophy of language that Tarski's formal result frames
but does not settle, and the various answers --- that natural language is
inconsistent in a harmless way, that it stratifies implicitly, that it
uses context to dodge the liar --- are live positions, not corollaries.
This book claims only the formal theorem: in formal languages of
sufficient strength, the truth predicate is undefinable, and truth is
recovered up the tower. Everything past the formal boundary is somebody's
philosophy, cited to them, never to Tarski.

The second, the non-funge note in its semantic form: the tower a given
object language sits under is \emph{its} tower. There is no universal
metalanguage crowning all languages at once --- that would be the very
top floor the theorem forbids --- so the ``truth'' one climbs to is always
truth-for-this-language, bound to the language below it, not a single
absolute Truth pooled across all houses. Each language's unnameable is its
own; the rooms, here as everywhere in this book, do not merge.

Two rooms now stand: G\"odel's inside the house, Tarski's cut into its
ascending wall, and they are visibly one construction aimed two ways. The
reader has surely begun to suspect that Cantor's, Russell's, G\"odel's,
and Tarski's rooms are not four discoveries but one, found four times. The
next chapter proves the suspicion --- not as a reading, this time, but as
a theorem with an author --- and that difference is the chapter's whole
point.
